% 题证摘录；来源与整理边界见相同编号分题文件。

\paragraph{Cartan's uniqueness theorem (support for the rigidity conclusion).}
\begin{thm}[Cartan]
Suppose $U\subset\C^n$ is a bounded domain, $a\in U$, $f:U\to U$ is a holomorphic mapping, $f(a)=a$, and $Df(a)$ is the identity. Then $f(z)=z$ for all $z\in U$.
\end{thm}
A polynomial $P:\C^n\to\C$ is homogeneous of degree $d$ if $P(sz)=s^dP(z)$ for all $s\in\C$ and $z\in\C^n$. A polynomial vector-valued mapping is homogeneous of degree $d$ if each component is. The power series at $a$ of a holomorphic $f$ is then written as $\sum_{m=0}^\infty f_m(z-a)$, where $f_m$ is a homogeneous polynomial of degree $m$.
\begin{proof}
Without loss of generality, assume $a=0$. Write $f$ as a power series at the origin in terms of homogeneous parts:
\[
f(z)=z+f_k(z)+\sum_{m=k+1}^\infty f_m(z)=z+f_k(z)+\text{higher-order terms},
\]
where $k\geq2$ is an integer such that $f_2(z)=f_3(z)=\cdots=f_{k-1}(z)=0$. The degree-one homogeneous part is simply the vector $z$, because the derivative of $f$ at the origin is the identity. Compose $f$ with itself $\ell$ times:
\[f^\ell(z)=\underbrace{f\circ f\circ\cdots\circ f}_{\ell\text{ times}}(z).\]
As $f(U)\subset U$, we have that $f^\ell$ is a holomorphic map of $U$ to $U$. As $U$ is bounded, there is an $M$ such that $\|z\|\leq M$ for all $z\in U$. Therefore, $\|f(z)\|\leq M$ for all $z\in U$, and $\|f^\ell(z)\|\leq M$ for all $z\in U$. Note that
\[f_k\bigl(f(z)\bigr)=f_k\bigl(z+\text{higher-order terms}\bigr)=f_k(z)+\text{higher-order terms}.\]
Therefore,
\begin{align*}
f^2(z)&=f\bigl(f(z)\bigr)=f(z)+f_k\bigl(f(z)\bigr)+\text{higher-order terms}\\
&=z+2f_k(z)+\text{higher-order terms}.
\end{align*}
Continuing this procedure,
\[f^\ell(z)=z+\ell f_k(z)+\text{higher-order terms}.\]
Suppose $\Delta_r(0)$ is a polydisc whose closure is in $U$. Via Cauchy estimates, for every multi-index $\alpha$ with $|\alpha|=k$,
\[
\frac{\alpha!}{r^\alpha}M\geq\left\|\frac{\partial^{|\alpha|}f^\ell}{\partial z^\alpha}(0)\right\|=\ell\left\|\frac{\partial^{|\alpha|}f}{\partial z^\alpha}(0)\right\|.
\]
The inequality holds for all $\ell\in\N$, and so $\frac{\partial^{|\alpha|}f}{\partial z^\alpha}(0)=0$. Therefore, $f_k\equiv0$. On the domain of convergence of the expansion, we get $f(z)=z$, as there is no other nonzero homogeneous part in the expansion of $f$. As $U$ is connected, the identity theorem says $f(z)=z$ for all $z\in U$.
\end{proof}

A circular domain is a domain $U\subset\C^n$ such that if $z\in U$, then $e^{i\theta}z\in U$ for all $\theta\in\R$.
\begin{cor}
Suppose $U,V\subset\C^n$ are bounded circular domains with $0\in U$, $0\in V$, and $f:U\to V$ is a biholomorphic map such that $f(0)=0$. Then $f$ is linear.
\end{cor}
\begin{proof}
The map $g(z)=f^{-1}\bigl(e^{-i\theta}f(e^{i\theta}z)\bigr)$ is an automorphism of $U$, and via the chain rule, $g'(0)=I$. Therefore, Cartan says that $f^{-1}\bigl(e^{-i\theta}f(e^{i\theta}z)\bigr)=z$, or in other words,
\[f(e^{i\theta}z)=e^{i\theta}f(z).\]
Write $f$ near zero as $f(z)=\sum_{m=1}^\infty f_m(z)$ where $f_m$ are homogeneous polynomials of degree $m$ (notice $f_0=0$). Then
\[
\sum_{m=1}^\infty e^{i\theta}f_m(z)=e^{i\theta}\sum_{m=1}^\infty f_m(z)=\sum_{m=1}^\infty f_m(e^{i\theta}z)=\sum_{m=1}^\infty e^{im\theta}f_m(z).
\]
By the uniqueness of the Taylor expansion, $e^{i\theta}f_m(z)=e^{im\theta}f_m(z)$, or $f_m(z)=e^{i(m-1)\theta}f_m(z)$, for all $m$, all $z$, and all $\theta$. If $m\neq1$, we obtain that $f_m\equiv0$, which proves the claim.
\end{proof}
